\maketitle
\begin{abstract}
We continue the polylogarithm manuscript at the pinned repository revision
identified in Appendix~\ref{app:provenance}. Four proof developments are obtained. First, an elementary
multiplicative-increment inequality controls every Gaussian Euler remainder
at arbitrary positive real orders. It identifies the least universal
constant as $C_*=\max_{b>0}\{\beta(b)+2^{-b}\eta(b)\}$, and exact rational
certificates enclose its unique maximizer and its value. Second, the
coordinate $y=\log\Gamma(x)$ gives an all-orders expansion of reflected
log-gamma moments uniformly for every ratio of their two nonnegative
exponents, as the larger exponent tends to infinity. Third, an explicit
Appell-root separation bound and a logarithmic variation estimate for
the Lerch kernel yield an effective saturation threshold and simultaneous
zero locations when the Laurent index grows with the derivative order.
Fourth, a central-factorial identity integrally diagonalizes the odd-weight
shuffle minor, giving its cokernel, Smith factors, field ranks, and optimal
common inverse denominator. Weight-eleven and weight-thirteen Gaussian
relations are recorded as conjectures with exact finite proximity
certificates. Established source results, new proofs, finite certificates,
numerical diagnostics, and unresolved claims are distinguished throughout.
\end{abstract}

\tableofcontents
\clearpage

\section{Scope, prior results, and the new claims}
\label{sec:introduction}

The starting point is the 375-page manuscript in
\code{Analysis/Polylogarithms/docs/manuscript}, together with the relevant
reports already present in the same repository at the pinned revision
\cite{repo}. The version matters: the $S_4$ identity, several Gaussian
reductions, and the rational relation ranks were already proved there.
Repeating those results as new conjecture resolutions would misstate the
state of the project. The existing $S_6$ special-value conjecture remains
unresolved by the present work.

The objective is to close specific gaps in the current programme and
provide reusable structural results. The universal Euler question has a
particularly direct answer: the existing Gaussian-axis maximum really is
the sharp error constant, once a new comparison is proved for \emph{all}
Euler indices. The reflected-moment question requires a different kind of
step. Separate fixed-exponent and bounded-ratio asymptotics do not justify
substitution of an unbounded exponent ratio. A change of coordinate and
an appropriate placement of the Jacobian in the phase give uniform
derivative and tail estimates on the full parameter range.

The zero problem has a similar obstruction to naive extrapolation. A
fixed-index expansion on a compact scale does not control the extreme
roots when the index grows. Our proof works with relative polynomial
errors on the logarithmic coordinate and with a scale-independent
positive-weight inequality. The resulting threshold is intentionally
conservative, but it is explicit, uniform in the Lerch parameter, and
applies to every positive root at once.

The algebraic part is separate from those analytic estimates. A rational
rank says that a coordinate system works over $\mathbb Q$. Integral
diagonalization additionally determines its torsion and all primes of
coordinate failure. The matrix is closely related to Zagier's double-zeta
matrix \cite{Zagier,Ma,MaTeo}; this connection is credited rather than
treated as evidence of new priority. The proof given here is constructive
and supplies the integral information directly in the manuscript's
shuffle coordinates.

\subsection{A claim and dependency map}

\begin{longtable}{@{}P{0.19\textwidth}P{0.46\textwidth}P{0.29\textwidth}@{}}
\toprule
Result & New conclusion & Inherited input or limitation\\
\midrule
\endhead
Euler remainder & The least constant uniform in $a,b>0$ and $N\ge1$ is
$C_*$; $C(b)$ is optimal at each fixed $b$. & The source's unique-axis-maximum
theorem identifies the sole maximizing $b_*$.\\
Exact maximum & Rational enclosures for $b_*$ and $C_*$, and the useful
budget $C_*<57/50$. & Finite exact arithmetic plus the proved qualitative
unimodality.\\
Reflected moments & All fixed asymptotic orders, uniformly for
$0\le m\le n$ as $n\to\infty$. & The previously certified global
log-gamma slope inequality; no exponentially small error is claimed.\\
Lerch zeros & Explicit $K_n$ and relative locations uniform in every
branch and $0\le\rho\le1$; a growing-index regime follows. & The
standard positive integral and Laplace zero-count principle, reviewed
in the proof. The threshold is sufficient, not sharp.\\
Shuffle lattice & Integral equivalence to
$\diag(1,3,\ldots,2m-1)$, hence the full arithmetic of the minor. & A
formal shuffle presentation; no period-independence conclusion.\\
Gaussian candidates & Frozen $S_{10}$ and $S_{12}$ formulas with exact
normalized residual bounds below $10^{-650}$. & They remain conjectures;
the residual intervals contain zero.\\
\bottomrule
\end{longtable}

All main analytic and algebraic statements are proved in ordinary
mathematical language. The exact certificates have explicitly narrower
roles: they certify constant brackets, check a pre-existing slope
inequality, or establish finite residual enclosures. They are not a
substitute for an identity proof, and no Lean formalization is claimed.

\subsection{Notation and branch conventions}

We use decreasing nested indices:
\[
 \Li_{a,b}(z,1)=\sum_{n>r\ge1}\frac{z^n}{n^a r^b},\qquad
 H_n^{(b)}=\sum_{r=1}^n r^{-b},\qquad H_n=H_n^{(1)}.
\]
The double family $F_{a,b}(z)$ always denotes this one-variable
specialization, initially in $|z|<1$. Real-order values at $i$ mean
the principal analytic continuation. At some small real orders the
ordinary boundary series does not converge, which does not affect the
finite Euler transforms or their proved limit. All gamma powers in real
integrals use positive real bases. The Dirichlet beta and eta functions
are denoted by $\beta$ and $\eta$; Euler's constant is $\gamma$.

In the moment section only, $F(x)=-\psi(x)$ is a positive real function;
the argument and absence of subscripts distinguish it from the double
polylogarithm. In the Lerch section $F_{n,k}^\rho$ is a signed normalization
of a parameter derivative, defined explicitly there. These local
notations follow the corresponding source problems and are not
identifications of different functions.

\subsection{Reading the computational evidence}

The accompanying archive contains the full LaTeX source, compiled PDF,
reusable section fragments, exact verifiers, frozen candidates, interval
certificates, numerical diagnostics, figure-generation code, a provenance
record, and integration notes. Each computational result says what it
establishes. For example, the equality $UMV=D$ is exact integer algebra;
an enclosure $|R|<10^{-650}$ is a finite analytic error statement; and
a plot of asymptotic errors is a floating-point diagnostic. None of these
labels is used interchangeably.

The original log-gamma moment problem and its fixed-order asymptotics
have substantial prior literature \cite{Amdeberhan,Bailey}. Our contribution
is the unrestricted-ratio theorem, rather than the definition or first
asymptotic observation. Likewise, the classical Laplace and gamma-series
tools used below are standard \cite{DLMFLaplace,DLMFGamma}; uniformity in
the present parameter domain is proved explicitly.
